\documentclass[12pt]{article}
\usepackage[T1]{fontenc}
\usepackage[utf8]{inputenc}
\usepackage{lmodern}
\usepackage{textcomp}
\usepackage{enumitem}
\usepackage{geometry}
\usepackage{graphicx}
\usepackage{amsmath, amssymb}
\geometry{margin=1in}
\begin{document}
\begin{center}
Chemistry - Graduate Level Assessment 24 \
Total Marks: 40 \
Answer all questions.
\end{center}

\section*{Multiple Choice (10 marks)}
\begin{enumerate}[label=\arabic*.]
\item A sample of solid MgCl 2 would be most soluble in which of the following solutions?
\begin{enumerate}[label=(\alph*)]
    \item AlCl 3(aq)
    \item CBr 4(aq)
    \item NaOH(aq)
    \item HCl(aq)
    \item LiOH(aq)
\end{enumerate}
\item 0.25 mol of a weak, monoprotic acid is dissolved in 0.25 L of distilled water. The pH was then measured as 4.26. What is the pKa of this weak acid?
\begin{enumerate}[label=(\alph*)]
    \item 4.26
    \item 3.66
    \item 9.26
    \item 2.66
    \item 10.52
\end{enumerate}
\item In the infrared spectrum of $\mathrm{H}^{127} \mathrm{I}$, there is an intense line at $2309 \mathrm{~cm}^{-1}$. Calculate the force constant of $\mathrm{H}^{127} \mathrm{I}$.
\begin{enumerate}[label=(\alph*)]
    \item 400 $ \mathrm{~N} \cdot \mathrm{m}^{-1}$
    \item 350 $ \mathrm{~N} \cdot \mathrm{m}^{-1}$
    \item 250 $ \mathrm{~N} \cdot \mathrm{m}^{-1}$
    \item 285 $\mathrm{~N} \cdot \mathrm{m}^{-1}$
    \item 450 $\mathrm{~N} \cdot \mathrm{m}^{-1}$
\end{enumerate}
\item Calculate the reduced mass of a nitrogen molecule in which both nitrogen atoms have an atomic mass of 14.00.
\begin{enumerate}[label=(\alph*)]
    \item 1.75
    \item 3.50
    \item 14.00
    \item 10.00
    \item 12.00
\end{enumerate}
\item Calculate the kinetic energy of an automobile weighing 2000 pounds and moving at a speed of 50 miles per hour. Record your answer in (a) Joules and (b) ergs.
\begin{enumerate}[label=(\alph*)]
    \item 3.0 \ensuremath{\times} $10^5$ Joules, 3.0 \ensuremath{\times} $10^12$ ergs
    \item 2.5 \ensuremath{\times} $10^5$ Joules, 2.5 \ensuremath{\times} $10^12$ ergs
    \item 2.0 \ensuremath{\times} $10^5$ Joules, 2.0 \ensuremath{\times} $10^12$ ergs
    \item 1.8 \ensuremath{\times} $10^5$ Joules, 1.8 \ensuremath{\times} $10^12$ ergs
    \item 3.56 \ensuremath{\times} $10^5$ Joules, 3.56 \ensuremath{\times} $10^12$ ergs
\end{enumerate}
\end{enumerate}

\section*{True / False (10 marks)}
\textit{Answer True or False}
\begin{enumerate}[label=\arabic*.]
\setcounter{enumi}{5}
\item True or False: The simplest alkene has a tetrahedral configuration around its double bond carbon atoms.
\item True or False: The number of allowed energy levels for the 55 Mn nuclide is 5.
\item True or False: The constants in the Dieterici equation are a' = 2 RT\_cV\_c and b' = V\_c / 2, derived from critical constants.
\item True or False: During the electrolysis of molten magnesium bromide, magnesium is produced at the cathode and bromine at the anode.
\item True or False: A light pulse of 2.00 mJ at a wavelength of 1.06 μm contains approximately 1.80 \ensuremath{\times} 10^\{16\} photons.
\end{enumerate}

\section*{Long Form Response (20 marks)}
\textit{Answer each question with a clear and complete explanation.}
\begin{enumerate}[label=\arabic*.]
\setcounter{enumi}{10}
\item Discuss the significance of the orbital angular momentum quantum number, l, in the context of ionization energy, particularly focusing on the electron that is most easily removed from ground-state aluminum.
\item Discuss how the decrease in atmospheric pressure at 4000 meters altitude affects the boiling point of water, and calculate the expected boiling point under these conditions.
\end{enumerate}

\vfill
\noindent\textit{End of Paper}
\end{document}